% 在 Overleaf 中使用 XeLaTeX；单文件、黑白排版。
% 数据来源：CPF_all_results完整归档；生成数据表见week5_evidence_20261005_verified。
\documentclass[12pt,a4paper,UTF8]{ctexart}

% 页面与基础宏包
\usepackage[top=2.5cm,bottom=2.5cm,left=2.7cm,right=2.7cm,headheight=15pt]{geometry}
\usepackage{array,tabularx,longtable,booktabs,enumitem,setspace,fancyhdr}
\usepackage{graphicx,float,caption,amsmath,amssymb,bm,titlesec,listings}

% 正文、表格、列表
\setstretch{1.5}
\setlength{\parindent}{2em}
\setlength{\parskip}{0pt}
\renewcommand{\arraystretch}{1.4}
\setlength{\tabcolsep}{5pt}
\setlength{\emergencystretch}{1.5em}
\clubpenalty=10000
\widowpenalty=10000
\displaywidowpenalty=10000
\captionsetup{font=small,labelfont=bf,labelsep=space}
\setlist[itemize]{leftmargin=2em,itemsep=0.2em,topsep=0.3em}
\setlist[enumerate]{leftmargin=2.5em,itemsep=0.2em,topsep=0.3em}

% 标题：黑体，全部黑白
\setcounter{secnumdepth}{4}
\setcounter{tocdepth}{4}
\titleformat{\paragraph}[block]{\normalfont\normalsize\bfseries}{\theparagraph}{1em}{}
\titleformat{\subparagraph}[block]{\normalfont\normalsize\bfseries}{\thesubparagraph}{1em}{}
\titlespacing*{\paragraph}{0pt}{1.5ex}{0.8ex}
\ctexset{
  section={format=\heiti\zihao{-3}\bfseries,name={,、},number=\chinese{section},beforeskip=1.5em,afterskip=0.8em},
  subsection={format=\heiti\zihao{4}\bfseries,beforeskip=1.0em,afterskip=0.5em},
  subsubsection={format=\heiti\zihao{-4}\bfseries,beforeskip=0.8em,afterskip=0.4em}
}
\renewcommand{\thesubsection}{\arabic{section}.\arabic{subsection}}
\renewcommand{\thesubsubsection}{\thesubsection.\arabic{subsubsection}}

% 页眉页脚
\newcommand{\reportauthor}{许瀚文}
\newcommand{\reportperiod}{第五周\quad 2026.9.28--2026.10.4}
\pagestyle{fancy}
\fancyhf{}
\fancyhead[R]{\zihao{5}\reportauthor}
\fancyfoot[C]{\thepage}
\renewcommand{\headrulewidth}{0.4pt}
\renewcommand{\footrulewidth}{0pt}

% 记录与表格命令
\newcommand{\fillin}[1][4cm]{\rule{#1}{0.4pt}}
\newcommand{\fillline}{\mbox{\fillin}\nobreak}
\newcommand{\file}[1]{\texttt{\detokenize{#1}}}
\newcolumntype{C}[1]{>{\centering\arraybackslash}p{#1}}
\newcolumntype{Y}{>{\raggedright\arraybackslash}X}
% Z 为自动分配宽度的居中列；避免与参考配置中的固定宽度 C{...} 冲突。
\newcolumntype{Z}{>{\centering\arraybackslash}X}
\newenvironment{reporttable}{\par\smallskip\begingroup\small\centering}{\par\endgroup\smallskip}

% 文件结构与 JSON 片段保持等宽字体；不使用彩色语法高亮。
\lstset{basicstyle=\ttfamily\small,breaklines=true,breakatwhitespace=false,
  columns=fullflexible,keepspaces=true,showstringspaces=false,tabsize=4,
  frame=single,framerule=0.3pt,aboveskip=0.45em,belowskip=0.5em,
  xleftmargin=1.5em,framexleftmargin=1em,xrightmargin=0pt}




\begin{document}
\begin{center}
{\zihao{2}\bfseries 科研组周报}\par
\vspace{0.6em}
{\normalsize\reportperiod}\par
\vspace{0.3em}
{\small 数据归档与核对截止：2026.10.5}
\end{center}
\vspace{0.5em}
\section{本周概述}
\subsection{概述}
本周先完成 CPF 在 UT-Zappos50K、MIT-States 和 C-GQA 上的基线复现与归档，再整理三个数据集的知识文件并构建 GloVe 缓存，随后完成 UT 的 A/C/D/E 对照、固定 checkpoint 的 $\lambda$ 后处理诊断，以及 MIT-States、C-GQA 的 A/C 扩展验证。报告依据本地完整归档中的原始指标、训练日志、有效配置和冻结清单填写。

主要发现是：UT 的知识文本版本 C 相对 A 获得较大验证集收益，CW/OW AUC 分别增加4.791/6.587个百分点；C-GQA 的对应增量仅为0.053/0.251个百分点；MIT-States 的对应增量为$-0.315/-0.010$个百分点。当前证据支持“文本表示在 UT 上有效、跨数据集效果存在差异”，尚不足以声称普遍稳定提升。UT 的固定 $K$ 校准不利于主要指标，Dpost/Epost 均选择 $\lambda=0$，保留为负结果。
\subsection{总体进展与完成边界}

\begin{reporttable}

\begin{tabularx}{\linewidth}{|p{3.4cm}|p{3.0cm}|Y|}
\hline

\textbf{工作} & \textbf{完成情况} & \textbf{证据与范围} \\ \hline

CPF 基线复现 & 已完成 & 3 个历史 run；test CW/OW 各1组 \\ \hline

知识整理与缓存 & 已完成结构核验 & 3 套 OI/OE/AE；冻结来源与缓存一致 \\ \hline

UT A/C/D/E & 已完成 & 4 个 run；每个10 epoch；8组 val 指标 \\ \hline

MIT/C-GQA A/C & 已完成 & 各2个 run；合计8组 val 指标 \\ \hline

$\lambda$ 后处理 & 已完成 & 2个固定 checkpoint；5个强度；20次评价 \\ \hline

条件知识增量/随机语义图 & 方案阶段 & EKR 与 USG 尚未训练，独立记录 \\ \hline

\end{tabularx}\end{reporttable}

历史基线共3个 run；本轮对照共8个 run，其中 A 的3个最佳权重与对应历史基线一致，不计为新的独立重复证据。所有本轮 run 的状态均为 \file{train_and_val_complete}，日志均包含第1至10轮验证。16份主指标与20份诊断指标的配置、checkpoint 和冻结清单核对通过；另外对11个归档 checkpoint 实际文件重新计算 SHA256，均与声明一致。

\noindent\textbf{统计口径：}指标按百分数显示，差值使用百分点（pp）；主表保留3位小数，全部精度另存 CSV。历史 test 与本轮 validation 分开分析。各数据集只有一个 seed；不填写不存在的均值、标准差或显著性结论。
\section{本周工作内容及总结}
\subsection{工作一：CPF 基线实验与复现结果}
\subsubsection{有效设置与版本}
使用 ViT-B/16 视觉模型与300维 GloVe，视觉输出维度为768，标准组合划分为 \file{compositional-split-natural}。三个基线实际均训练10轮，训练/评价 batch size 均为64，workers为4；每轮按 validation CW AUC 选择最佳 checkpoint，再对该权重进行 test CW/OW 评价。GloVe 原语嵌入在当前代码中冻结，视觉 backbone 微调。

模型参数学习率为 $10^{-4}$，视觉编码器为 $10^{-5}$，名称嵌入配置项为 $2.5\times10^{-6}$，权重衰减为 $5\times10^{-5}$；冻结的嵌入不因配置学习率而产生更新。cosine 分类温度为0.05。配置中的衰减节点23/30和增强起始轮37均未在本轮10轮预算内触发。

\begin{reporttable}

\begin{tabularx}{\linewidth}{|Y|r|r|r|c|}
\hline

\textbf{数据集} & \textbf{seed} & \textbf{实际轮数} & \textbf{最佳轮} & \textbf{训练/评价批次} \\ \hline

UT-Zappos & 3351 & 10 & 6 & 64/64 \\ \hline

MIT-States & 0 & 10 & 10 & 64/64 \\ \hline

C-GQA & 3351 & 10 & 9 & 64/64 \\ \hline

\end{tabularx}\end{reporttable}

服务器记录 Python 3.8.18，PyTorch版本为1.13.1+cu117，torchvision版本为0.14.1+cu117；CUDA runtime为11.7，GPU为 RTX 3090。当前 pip 清单还记录 NumPy 1.24.4、Pillow 10.4.0。历史代码 commit 为：
\par\noindent{\small\file{690467371295943f3f49cb9e856bb8401bac7c03}}\par
存在本地改动，故同时保存代码 patch、源码快照和文件清单，而不只用 commit 代表运行代码。

发现归档的部分 \file{config_used.yml} 保留了20轮、旧路径或 seed=-1 等基础配置，未包含命令行覆盖项。本报告采用 \file{train.log} 开头打印的有效配置：MIT 的实际 seed=0、batch=64，MIT/C-GQA 的实际预算均为10轮。该差异已单独记录，后续统一保存合并后的训练配置。
\subsubsection{已归档 test 结果}

\begin{reporttable}

\begin{tabularx}{\linewidth}{|Y|c|r|r|r|r|}
\hline

\textbf{数据集} & \textbf{设置} & \textbf{AUC\%} & \textbf{HM\%} & \textbf{Seen\%} & \textbf{Unseen\%} \\ \hline

UT-Zappos & CW & 36.51 & 52.59 & 61.49 & 68.48 \\ \hline

UT-Zappos & OW & 27.07 & 43.87 & 61.49 & 51.30 \\ \hline

MIT-States & CW & 11.05 & 26.57 & 39.66 & 35.81 \\ \hline

MIT-States & OW & 4.18 & 14.67 & 39.66 & 14.38 \\ \hline

C-GQA & CW & 6.81 & 21.30 & 38.42 & 20.37 \\ \hline

C-GQA & OW & 1.58 & 7.27 & 38.42 & 4.98 \\ \hline

\end{tabularx}\end{reporttable}

\noindent{\small 来源：3个历史 run 的 \file{repro_meta/final_results.txt} 与训练/评价日志。Seen/Unseen 是扫描校准曲线后各自的最佳值，不一定处于同一个偏置点，不能直接代入调和均值重算 best-HM。\par}

\subsubsection{与论文参考及复现问题}

\begin{reporttable}

\begin{tabularx}{\linewidth}{|Y|r|r|r|r|}
\hline

\textbf{数据集/设置} & \textbf{论文AUC\%} & \textbf{论文HM\%} & \textbf{$\Delta$AUC(pp)} & \textbf{$\Delta$HM(pp)} \\ \hline

UT-Zappos/CW & 41.40 & 55.70 & -4.89 & -3.11 \\ \hline

UT-Zappos/OW & 31.20 & 47.60 & -4.13 & -3.73 \\ \hline

MIT-States/CW & 11.20 & 26.80 & -0.15 & -0.23 \\ \hline

MIT-States/OW & 4.40 & 15.10 & -0.22 & -0.43 \\ \hline

C-GQA/CW & 8.20 & 23.90 & -1.39 & -2.60 \\ \hline

C-GQA/OW & 2.10 & 9.50 & -0.52 & -2.23 \\ \hline

\end{tabularx}\end{reporttable}

\noindent{\small 论文参考为 CPF（ICCV 2025）本地论文第6页表1、表2的 ViT-B 行，已核对原文；差值为本次复现减论文报告值，不代表新增方法的收益。\par}

MIT-States 最接近论文报告值；UT 与 C-GQA 的 AUC/HM 仍有差距，当前为部分复现。论文正文同样描述10轮训练，故不能仅以“训练轮数不足”解释差距。还需逐项检查视觉预训练权重、图像预处理、数据标签集合、评分分支、数值精度和随机种子。现有记录只能定位核对项，不能确定某一项已构成误差根因。

CW 到 OW 后，三个基线的 AUC/HM 均下降。C-GQA 历史 test OW AUC 为1.58，而本轮 A 的 val OW AUC 为7.574；两者对应不同数据阶段及未见组合集合，不能相减后宣称方法涨点。三个当前 A 的最佳权重哈希分别与历史基线相同，这进一步说明必须区分数据阶段。
\subsection{工作二：知识整理、上传与结构记录}
\subsubsection{上传目录与字段含义}
服务器在 \file{knowledge/kcc0601_v1} 下保存知识文件。实际训练来源子目录为 \file{trusted_v1}，缓存子目录为 \file{cache_trusted_v1}。\file{draft} 保留重构草稿，\file{reviewed_v1} 是另一目录，不能仅凭目录名推断它就是本轮使用版本。训练名称 \file{ut-zappos} 对应知识和数据目录 \file{ut-zap50k}。

\par\vspace{0.5\baselineskip}\noindent\begin{minipage}{\linewidth}

\begin{lstlisting}

trusted_v1/
  ut-zap50k/  mit-states/  cgqa/
    attributes.txt       objects.txt
    oi_dict.json         oe_dict.json         ae_dict.json
    review_status.csv    entry_provenance.csv
    manifest.json
    K_oa_text_lexical.csv K_oa_text_report.json
cache_trusted_v1/<dataset>/
    W_o_K.pt             W_a_K.pt
    W_o_name_raw.pt      W_a_name_raw.pt
    W_o_name_unit.pt     W_a_name_unit.pt
    K_oa_text_raw.pt     K_oa_text_template.pt
    K_oa_text.pt         K_oa_prior_centered.pt
    W_ao_K/obj_XXXXX.pt  labels.json
    manifest.json       cache_manifest.json   report.json

\end{lstlisting}\end{minipage}\par\smallskip

\begin{reporttable}

\begin{tabularx}{\linewidth}{|p{5cm}|Y|}
\hline

\textbf{文件} & \textbf{结构与作用} \\ \hline

\file{oi_dict.json} & 父类名称→对象标签列表；一个对象仅归属一个父类 \\ \hline

\file{oe_dict.json} & 对象标签→3条描述文本 \\ \hline

\file{ae_dict.json} & 属性标签→3条描述文本 \\ \hline

\file{attributes/objects.txt} & 原语标签及固定行序，与训练 reader 和缓存一致 \\ \hline

\file{review_status.csv} & kind、label、status、reviewer、date、note \\ \hline

\file{entry_provenance.csv} & query、synset、sense\_count、head fallback、parent、quality、review reason \\ \hline

\file{manifest.json} & 来源、状态、数量、许可/非官方说明及知识文件哈希 \\ \hline

\end{tabularx}\end{reporttable}

UT 的实际条目示例如下，标签保持数据集原始写法：

\par\vspace{0.5\baselineskip}\noindent\begin{minipage}{\linewidth}

\begin{lstlisting}

OI: {"boot": ["Boots.Ankle", "Boots.Knee.High",
              "Boots.Mid-Calf"]}
OE: {"Boots.Ankle": ["shaft ending at the ankle",
      "closed toe boot silhouette", "low boot collar"]}
AE: {"Canvas": ["woven canvas fabric",
      "coarse textile surface", "matte cloth finish"]}

\end{lstlisting}\end{minipage}\par\smallskip

\subsubsection{数量、来源与语义审核状态}

\begin{reporttable}

\begin{tabularx}{\linewidth}{|Y|r|r|r|r|r|r|}
\hline

\textbf{数据集} & \textbf{属性} & \textbf{对象} & \textbf{父类} & \textbf{OE描述} & \textbf{AE描述} & \textbf{审核行数} \\ \hline

UT-Zappos & 16 & 12 & 4 & 36 & 48 & 28 \\ \hline

MIT-States & 115 & 245 & 19 & 735 & 345 & 360 \\ \hline

C-GQA & 413 & 674 & 23 & 2022 & 1239 & 1087 \\ \hline

\end{tabularx}\end{reporttable}

UT 为先前手工重构的知识文本；MIT/C-GQA 为 Princeton WordNet 3.0 首义项词典草稿。三个版本均不是 LogiCZSL 官方知识文件。来源记录中，UT 有28条 \file{manual_reconstruction_unreviewed}；MIT 有360条 \file{wordnet_auto_unreviewed}；C-GQA 有1081条自动 WordNet 记录和6条 \file{label_template_only}。

审核表中的 approved 行数分别为28、360、1087，pending为0，但冻结 manifest 仍标记 \file{bulk_trusted_unverified}，并明确不与正式审核结果等价。这些条目采用了批量信任流程；结构验证、审核状态修改和低 OOV 都不证明语义逐条正确。尤其多义词首义项、仅有名称模板的条目和非视觉定义，需要后续抽样复核。

本次对每个数据集的7个必需知识源文件重新计算哈希，与缓存 \file{input_sha256} 一致；标签内容与顺序一致。来源 CSV 是补充材料，未作为独立输入冻结进缓存，其信息不能替代已核验的 OI/OE/AE 原文。
\subsubsection{GloVe 编码与模型实际使用方式}

\begin{reporttable}

\begin{tabularx}{\linewidth}{|Y|c|c|c|r|}
\hline

\textbf{数据集} & \textbf{$W_o^K$形状} & \textbf{$W_a^K$形状} & \textbf{$K$形状} & \textbf{OOV比例} \\ \hline

UT-Zappos & 12×300 & 16×300 & 12×16 & 0.361\% \\ \hline

MIT-States & 245×300 & 115×300 & 245×115 & 0.260\% \\ \hline

C-GQA & 674×300 & 413×300 & 674×413 & 0.337\% \\ \hline

\end{tabularx}\end{reporttable}

名称词向量复用基线别名与同一 GloVe 词表；描述经 token 化后，先平均词表内 token 向量，再逐句 L2 归一化，最后将名称单位向量与3条描述单位向量平均并归一化：
\[
W_i^K=\operatorname{norm}\!\left(\frac{1}{4}\sum_{r=0}^{3}t_{ir}\right),
\qquad t_{i0}=\operatorname{norm}(W_i^{name}).
\]
三数据集描述的 fallback sentence 均为0。本次直接审计18个主要缓存张量的实际形状、有限值和哈希，知识原语向量的单位范数检查通过，中心化先验与保存的 $K$ 一致。

C/E 在训练与推理中使用冻结的 $W_o^K,W_a^K$；D/E 使用中心化 $K$ 进行分数校准。OI 父类结构和 \file{W_ao_K} 的组合文本块未进入本轮 MVP 图传播，也没有训练逻辑模块。因而本周实验是文本表示与静态校准验证，不能写成随机语义图或完整逻辑知识模型已经实现。
\subsubsection{组合先验构建与分布检查}
对全部对象×属性组合使用相同三种模板，包括“a object that is attribute”、“the attribute object”和“attribute appearance of object”，其中 object/attribute以标签替换。组合模板向量与 $\operatorname{norm}(W_o^K+W_a^K)$ 的 cosine 经全矩阵 min--max 得到 $K\in[0,1]$。选用的是 template 版本；目录中的 lexical CSV 仅为另一候选，未作为本轮 $K$。
\[
\widetilde K=\frac{K-\operatorname{mean}(K)}{\operatorname{std}_{pop}(K)},
\quad s'=s+\lambda\widetilde K_{oa}.
\]
标准差使用总体定义；常数 $K$ 的先验为零。当前 $K$ 只代表文本相似度，不是经校验的组合可行性或 LLM 事实评分。

\begin{reporttable}

\footnotesize

\begin{tabularx}{\linewidth}{|Y|r|r|r|r|r|}
\hline

\textbf{数据集} & \textbf{$K$均值} & \textbf{$K$标准差} & \textbf{中位数} & \textbf{零值比例} & \textbf{原始标准差} \\ \hline

UT-Zappos & 0.508405 & 0.244008 & 0.515478 & 0.520833\% & 0.041497 \\ \hline

MIT-States & 0.722081 & 0.120815 & 0.739266 & 0.003549\% & 0.021287 \\ \hline

C-GQA & 0.756356 & 0.092073 & 0.768898 & 0.000359\% & 0.021576 \\ \hline

\end{tabularx}\end{reporttable}

三组归一化 $K$ 的 min/max 均为0/1，中心化后全矩阵均值近0、总体标准差为1；非有限值、常数退化及均匀先验恒等检查均通过。随后只读取 train/val pair 列表，比较中心化先验的逐组合分布：

\begin{reporttable}

\footnotesize

\begin{tabularx}{\linewidth}{|Y|c|r|r|r|r|r|}
\hline

\textbf{数据集} & \textbf{pair集合} & \textbf{数量} & \textbf{均值} & \textbf{总体标准差} & \textbf{min} & \textbf{max} \\ \hline

UT-Zappos & train & 83 & 0.182 & 0.950 & -1.638 & 2.015 \\ \hline

UT-Zappos & val & 30 & -0.124 & 0.900 & -2.030 & 1.838 \\ \hline

MIT-States & train & 1262 & 0.131 & 0.975 & -4.345 & 2.300 \\ \hline

MIT-States & val & 600 & 0.172 & 0.924 & -3.611 & 2.300 \\ \hline

C-GQA & train & 5592 & 0.639 & 0.713 & -4.548 & 2.450 \\ \hline

C-GQA & val & 2292 & 0.746 & 0.672 & -3.668 & 2.359 \\ \hline

\end{tabularx}\end{reporttable}

UT 的 train 均值0.182、val均值$-0.124$，表明该 $K$ 在组合集合层面对训练组合更有利；这与正强度校准降低未见指标的结果相容，但不能仅凭分布差异证明逐图错误的因果机制。MIT/C-GQA 的 val均值并未低于 train，因此也不能把 UT 的解释直接推广到所有数据集。这里每个 pair 等权，未按图像数量加权。

语义抽查中，UT 的高分包括“Full.grain.leather--Boots.Knee.High”，低分包括“Rubber--Shoes.Oxfords”；C-GQA 的最高分为“used--person”，还出现“paper--person”等高分组合。这提示模板相似度可能反映词义、常用词或共享句式，不能自动解释成视觉兼容性。后续审核应保存判定依据，不能按 test 效果反向修改 $K$。
\subsection{工作三：ACDE 对照、校准诊断与跨数据集验证}
\subsubsection{对照定义与共同协议}

\begin{reporttable}

\begin{tabularx}{\linewidth}{|c|Y|c|c|}
\hline

\textbf{变体} & \textbf{原语文本} & \textbf{$K$校准} & \textbf{强度} \\ \hline

A & 基线名称向量 & 关闭 & 0 \\ \hline

C & 名称与3条描述聚合 & 关闭 & 0 \\ \hline

D & 基线名称向量 & 开启 & 0.1 \\ \hline

E & 名称与3条描述聚合 & 开启 & 0.1 \\ \hline

\end{tabularx}\end{reporttable}

全部新 run 使用相同10轮预算、有效 batch=64、相同损失权重、视觉 backbone、数据划分和 \file{legacy} 评分方式；逻辑/特征模块关闭。UT seed=3351，MIT=0，C-GQA=3351。各数据集 A/C 除文本开关及输出目录等预期项外，有效配置匹配；代码清单、知识缓存、数据清单、候选顺序和环境声明一致。

每个 run 以 val CW AUC 选最佳权重，同一权重再评价 val CW 与 OW。\file{kcc_test_at_end=False}，本轮没有生成新方法的 test 结果；训练末尾日志中的 test=0 是未执行测试的占位输出，不是模型测试准确率为零。

\begin{reporttable}

\footnotesize

\begin{tabularx}{\linewidth}{|Y|r|r|r|r|r|r|}
\hline

\textbf{数据集} & \textbf{train图像} & \textbf{val图像} & \textbf{train pair} & \textbf{val pair} & \textbf{CW候选} & \textbf{OW候选} \\ \hline

UT-Zappos & 22998 & 3214 & 83 & 30 & 116 & 192 \\ \hline

MIT-States & 30338 & 10420 & 1262 & 600 & 1962 & 28175 \\ \hline

C-GQA & 26920 & 7280 & 5592 & 2292 & 7767 & 278362 \\ \hline

\end{tabularx}\end{reporttable}

三数据集 smoke 均通过：A 与原基线分数最大差为0且 top-1一致；$\lambda=0$ 恒等、均匀先验恒等、候选重排与先验轴顺序正确；训练梯度有限且非零，知识缓存保持冻结。这些检查验证接入实现，不替代性能或知识语义验证。

\subsubsection{UT-Zappos 的 A/C/D/E 完整验证结果}

\noindent\textbf{Validation CW}\par\nopagebreak[4]

\begin{reporttable}

\begin{tabularx}{\linewidth}{|Z|Z|Z|Z|Z|Z|}
\hline

\textbf{变体} & \textbf{AUC\%} & \textbf{HM\%} & \textbf{Seen\%} & \textbf{Unseen\%} & \textbf{$\Delta$AUC(pp)} \\ \hline

A & 39.470 & 54.841 & 67.959 & 66.410 & +0.000 \\ \hline

C & 44.262 & 57.894 & 68.871 & 71.502 & +4.791 \\ \hline

D & 32.464 & 51.795 & 65.564 & 56.611 & -7.006 \\ \hline

E & 37.239 & 53.897 & 66.135 & 64.741 & -2.231 \\ \hline

\end{tabularx}\end{reporttable}

\noindent\textbf{Validation OW}\par\nopagebreak[4]

\begin{reporttable}

\begin{tabularx}{\linewidth}{|Z|Z|Z|Z|Z|Z|}
\hline

\textbf{变体} & \textbf{AUC\%} & \textbf{HM\%} & \textbf{Seen\%} & \textbf{Unseen\%} & \textbf{$\Delta$AUC(pp)} \\ \hline

A & 25.492 & 37.269 & 67.959 & 53.487 & +0.000 \\ \hline

C & 32.079 & 46.367 & 68.871 & 59.435 & +6.587 \\ \hline

D & 19.107 & 31.700 & 65.564 & 45.272 & -6.385 \\ \hline

E & 23.972 & 36.262 & 66.135 & 53.830 & -1.520 \\ \hline

\end{tabularx}\end{reporttable}

\noindent{\small 来源为原始 metrics.json；UT D 的 CW Unseen 实际为56.611\%，已按原始精度核对。\par}

\subsubsection{文本与校准的比较关系}

\begin{reporttable}

\begin{tabularx}{\linewidth}{|Z|Z|Z|Z|Z|}
\hline

\textbf{比较} & \textbf{CW $\Delta$AUC} & \textbf{CW $\Delta$HM} & \textbf{OW $\Delta$AUC} & \textbf{OW $\Delta$HM} \\ \hline

C-A & +4.791 & +3.053 & +6.587 & +9.097 \\ \hline

D-A & -7.006 & -3.046 & -6.385 & -5.569 \\ \hline

E-C & -7.022 & -3.997 & -8.107 & -10.105 \\ \hline

E-D & +4.775 & +2.101 & +4.865 & +4.562 \\ \hline

E-A & -2.231 & -0.944 & -1.520 & -1.008 \\ \hline

E-C-D+A & -0.016 & -0.952 & -1.722 & -4.536 \\ \hline

\end{tabularx}\end{reporttable}

C-A 是当前文本表示整体改动的效果：CW/OW AUC增加4.791/6.587 pp，HM增加3.053/9.097 pp。D-A 显示仅加固定校准后 AUC下降7.006/6.385 pp；E-C 显示在文本基础上加入校准后 AUC下降7.022/8.107 pp。E 仍低于 A，说明文本改善未抵消当前校准带来的损失。

E-D 仍为正值，说明在有校准的背景下，文本表示也能改善结果。交互差 $E-C-D+A$ 的 AUC为$-0.016/-1.722$ pp；CW 近零，OW为负。本表是单次训练的描述性差值，不是统计检验或因果效应分解。

需要区分原语向量变化与评价曲线变化。例如 UT CW 的 best-HM 点，A 的 $(S,U)$ 为$(50.969,59.350)\%$，C 为$(55.644,60.334)\%$；OW 中对应为$(40.479,34.531)\%$、$(52.794,41.335)\%$。这些点可重算相应 HM；它们与表中的各自最佳 Seen/Unseen 不是同一组数值。
\subsubsection{固定 checkpoint 的 $\lambda$ 后处理诊断}
Dpost 在 A 最佳权重上校准，Epost 在 C 最佳权重上校准；每条路线只改变推理 $\lambda$，不重新训练。网格为 $\{0,0.025,0.05,0.1,0.2\}$，每个强度分别评估 CW/OW。它们与重训 D/E 是两种不同实验。

\noindent\textbf{Dpost（A权重）}

\begin{reporttable}

\footnotesize

\begin{tabularx}{\linewidth}{|Z|Z|Z|Z|Z|Z|Z|}
\hline

\textbf{$\lambda$} & \textbf{CW AUC\%} & \textbf{CW HM\%} & \textbf{OW AUC\%} & \textbf{OW HM\%} & \textbf{OW U\%} & \textbf{入选资格} \\ \hline

0.000 & 39.470 & 54.841 & 25.492 & 37.269 & 53.487 & 通过 \\ \hline

0.025 & 38.399 & 55.024 & 24.449 & 36.775 & 52.931 & 未通过 \\ \hline

0.050 & 36.775 & 54.211 & 22.719 & 35.208 & 50.492 & 未通过 \\ \hline

0.100 & 31.572 & 51.037 & 18.348 & 30.782 & 43.902 & 未通过 \\ \hline

0.200 & 20.402 & 40.479 & 11.897 & 25.168 & 33.248 & 未通过 \\ \hline

\end{tabularx}\end{reporttable}

\noindent\textbf{Epost（C权重）}

\begin{reporttable}

\footnotesize

\begin{tabularx}{\linewidth}{|Z|Z|Z|Z|Z|Z|Z|}
\hline

\textbf{$\lambda$} & \textbf{CW AUC\%} & \textbf{CW HM\%} & \textbf{OW AUC\%} & \textbf{OW HM\%} & \textbf{OW U\%} & \textbf{入选资格} \\ \hline

0.000 & 44.262 & 57.894 & 32.079 & 46.367 & 59.435 & 通过 \\ \hline

0.025 & 43.963 & 57.209 & 31.882 & 45.589 & 60.419 & 未通过 \\ \hline

0.050 & 42.577 & 56.584 & 30.529 & 43.931 & 59.949 & 未通过 \\ \hline

0.100 & 38.187 & 54.816 & 25.842 & 39.163 & 55.413 & 未通过 \\ \hline

0.200 & 26.807 & 46.666 & 16.993 & 31.480 & 43.603 & 未通过 \\ \hline

\end{tabularx}\end{reporttable}

选择规则为在合格候选中最大化 val CW AUC，并相对各自锚点约束：$\Delta$CW HM$\ge-0.5$ pp、$\Delta$OW AUC$\ge-0.5$ pp、$\Delta$OW Unseen$\ge-2$ pp。两条路线的 $\lambda=0$ 均与锚点指标完全相同，且均为最终入选值；所有正强度候选都未通过保护约束，也都未超过锚点 CW AUC。

小强度并非所有单项都下降：Dpost 0.025 的 CW HM增加0.183 pp，但 OW AUC下降1.044 pp，故不合格；Epost 0.025 的 OW Unseen增加0.984 pp，但 CW HM下降0.686 pp，亦不合格。由此选择零意味着当前知识版本、checkpoint、网格和规则下不采用校准，不能扩展为“全部外部知识无效”。继续扩大固定 $K$ 强度不作为当前主路线。
\subsubsection{MIT-States 与 C-GQA 的 A/C 验证结果}

\begin{reporttable}

\begin{tabularx}{\linewidth}{|Y|c|r|r|r|r|}
\hline

\textbf{数据集} & \textbf{变体/设置} & \textbf{AUC\%} & \textbf{HM\%} & \textbf{Seen\%} & \textbf{Unseen\%} \\ \hline

MIT-States & A/CW & 13.707 & 30.227 & 41.594 & 40.403 \\ \hline

MIT-States & C/CW & 13.392 & 30.009 & 40.672 & 40.205 \\ \hline

MIT-States & A/OW & 4.602 & 15.754 & 41.594 & 14.715 \\ \hline

MIT-States & C/OW & 4.592 & 15.831 & 40.672 & 14.937 \\ \hline

C-GQA & A/CW & 8.689 & 23.929 & 37.649 & 28.094 \\ \hline

C-GQA & C/CW & 8.742 & 24.314 & 37.568 & 28.149 \\ \hline

C-GQA & A/OW & 7.574 & 22.325 & 37.649 & 24.582 \\ \hline

C-GQA & C/OW & 7.825 & 23.002 & 37.568 & 24.582 \\ \hline

\end{tabularx}\end{reporttable}

\noindent\textbf{三数据集统一比较：C减A（单位为pp）}

\begin{reporttable}

\begin{tabularx}{\linewidth}{|Y|c|r|r|r|r|}
\hline

\textbf{数据集} & \textbf{设置} & \textbf{$\Delta$AUC} & \textbf{$\Delta$HM} & \textbf{$\Delta$Seen} & \textbf{$\Delta$Unseen} \\ \hline

UT-Zappos & CW & +4.791 & +3.053 & +0.912 & +5.092 \\ \hline

UT-Zappos & OW & +6.587 & +9.097 & +0.912 & +5.948 \\ \hline

MIT-States & CW & -0.315 & -0.217 & -0.922 & -0.198 \\ \hline

MIT-States & OW & -0.010 & +0.077 & -0.922 & +0.222 \\ \hline

C-GQA & CW & +0.053 & +0.385 & -0.081 & +0.056 \\ \hline

C-GQA & OW & +0.251 & +0.677 & -0.081 & +0.000 \\ \hline

\end{tabularx}\end{reporttable}

UT 的方向和幅度最明确；C-GQA 六轮后逐渐追近基线，但最终 CW AUC只增加0.053 pp，不能据单 seed 小差值声称可靠提升。MIT 的 CW AUC/HM略降，OW AUC近乎持平，OW HM增加0.077 pp；这是混合结果，应保留为效果边界。

按此前操作方案的投入门槛，三个数据集中有两个 CW AUC为正，且 OW 未出现明显崩塌，满足继续做小规模重复验证的工程条件。但该门槛只决定是否继续投入，不是统计显著性标准。现阶段不能用 UT 的较大增量掩盖 MIT 的负结果，也不能将三个数据集简单合并成一个“平均涨点”结论。
\subsubsection{训练收敛与运行开销}

\begin{reporttable}

\footnotesize

\begin{tabularx}{\linewidth}{|Y|r|r|r|r|r|}
\hline

\textbf{run} & \textbf{最佳轮} & \textbf{总时长(min)} & \textbf{训练峰值} & \textbf{OW时长(min)} & \textbf{OW峰值} \\ \hline

UT-Zappos/A & 6 & 82.2 & 12296.1 & 0.49 & 2395.2 \\ \hline

UT-Zappos/C & 6 & 77.6 & 12296.1 & 0.38 & 2395.3 \\ \hline

UT-Zappos/D & 6 & 74.5 & 12296.1 & 0.40 & 2395.2 \\ \hline

UT-Zappos/E & 6 & 73.8 & 12296.1 & 0.37 & 2395.3 \\ \hline

MIT-States/A & 10 & 74.7 & 12318.0 & 4.91 & 2428.5 \\ \hline

MIT-States/C & 10 & 69.7 & 12318.4 & 4.91 & 2428.9 \\ \hline

C-GQA/A & 9 & 154.8 & 12386.4 & 30.57 & 5000.0 \\ \hline

C-GQA/C & 10 & 158.9 & 12387.6 & 31.31 & 5001.2 \\ \hline

\end{tabularx}\end{reporttable}

\noindent{\small 总时长为 \file{end_time-start_time}，含训练、逐轮验证和最终 CW/OW 评价，不是纯训练时间。峰值单位实际为 MiB，均为当前进程 PyTorch allocated memory，不等于整卡占用；运行时长受并行负载影响，不能由此声称 C 更快。\par}

八个 run 的逐轮 val CW AUC如下，数值来自训练日志，保留日志打印精度；最终主结果来自最佳权重的原始 JSON：

\begin{reporttable}

\scriptsize

\begin{tabularx}{\linewidth}{|Z|Z|Z|Z|Z|Z|Z|Z|Z|}
\hline

\textbf{轮} & \textbf{UT-A} & \textbf{UT-C} & \textbf{UT-D} & \textbf{UT-E} & \textbf{MIT-A} & \textbf{MIT-C} & \textbf{CG-A} & \textbf{CG-C} \\ \hline

1 & 28.23 & 37.04 & 18.53 & 23.39 & 5.01 & 4.67 & 4.60 & 4.33 \\ \hline

2 & 36.92 & 41.16 & 26.79 & 34.79 & 8.42 & 7.69 & 6.02 & 5.58 \\ \hline

3 & 38.89 & 43.32 & 28.98 & 35.83 & 9.64 & 9.08 & 6.85 & 6.84 \\ \hline

4 & 37.27 & 42.01 & 29.54 & 36.04 & 10.67 & 10.09 & 7.43 & 7.51 \\ \hline

5 & 39.11 & 43.85 & 30.69 & 36.26 & 11.79 & 11.05 & 7.95 & 8.03 \\ \hline

6 & 39.47 & 44.26 & 32.46 & 37.24 & 12.57 & 11.31 & 8.16 & 8.10 \\ \hline

7 & 36.63 & 41.29 & 30.47 & 36.11 & 12.71 & 11.86 & 8.33 & 8.35 \\ \hline

8 & 38.22 & 42.34 & 32.39 & 36.83 & 12.98 & 12.30 & 8.21 & 8.36 \\ \hline

9 & 34.27 & 38.82 & 27.51 & 32.67 & 12.97 & 12.45 & 8.69 & 8.72 \\ \hline

10 & 37.66 & 40.88 & 32.02 & 34.34 & 13.71 & 13.39 & 8.51 & 8.74 \\ \hline

\end{tabularx}\end{reporttable}

UT 四个变体最佳轮均为6；此后 loss仍下降，但 A/C 的 val AUC分别由39.47/44.26降至末轮37.66/40.88，表现为验证收益未随训练 loss持续改善。MIT A/C最佳轮均为10，末轮仍在提高；C-GQA A最佳轮为9、C为10。预算末端最优提示收敛可能尚未充分比较，不能证明加长训练一定有益；如延长，必须对 A/C 使用相同预算，独立命名并保留本轮结果。

C-GQA OW有278362个候选、7280张图像，最终一次评价约30.57/31.31分钟，显著长于 UT 的不足半分钟。这是候选规模和评价计算导致的开销，训练完成后继续运行 OW 评价属于正常阶段。
\subsubsection{保存预测的逐图误差分析}
按 \file{labels.npz} 中保存的 $bias=0$ top-1，核对 A/C 的图像顺序、候选顺序和真值一致后逐图配对。seen 身份仅由训练 pair 列表判定；不以 test 结果或 test pair 列表选择策略。以下为 OW 的原始分数 top-1 结果：

\begin{reporttable}

\footnotesize

\begin{tabularx}{\linewidth}{|Y|c|r|r|r|r|r|}
\hline

\textbf{数据集} & \textbf{分组} & \textbf{图像数} & \textbf{A准确率\%} & \textbf{C准确率\%} & \textbf{$\Delta$(pp)} & \textbf{净纠正数} \\ \hline

UT-Zappos & 全体 & 3214 & 21.811 & 21.935 & +0.124 & 4 \\ \hline

UT-Zappos & seen & 877 & 67.503 & 68.301 & +0.798 & 7 \\ \hline

UT-Zappos & val unseen & 2337 & 4.664 & 4.536 & -0.128 & -3 \\ \hline

MIT-States & 全体 & 10420 & 11.334 & 11.267 & -0.067 & -7 \\ \hline

MIT-States & seen & 1844 & 36.497 & 35.412 & -1.085 & -20 \\ \hline

MIT-States & val unseen & 8576 & 5.924 & 6.075 & +0.152 & 13 \\ \hline

C-GQA & 全体 & 7280 & 22.665 & 23.475 & +0.810 & 59 \\ \hline

C-GQA & seen & 3692 & 35.699 & 35.672 & -0.027 & -1 \\ \hline

C-GQA & val unseen & 3588 & 9.253 & 10.925 & +1.672 & 60 \\ \hline

\end{tabularx}\end{reporttable}

UT 全体净纠正4张（43张由错变对、39张由对变错），但 val unseen净减少3张；它的 AUC/HM明显改善主要体现在扫描偏置后的评价曲线，而不表现为原始 $bias=0$ 未见准确率普遍改善。MIT 全体净减少7张，seen减少20张、val unseen增加13张；说明不同组之间存在取舍。C-GQA 全体净纠正59张，其中 val unseen净纠正60张、seen净减少1张，未见组原始 top-1 增加1.672 pp。

这些计数描述同一批验证图像上的预测变化，不构成新 seed 的独立证据。原始 top-1 也不是 evaluator 的 best-HM 点或其固定高偏置输出，因此不与 \file{closed_match}、best Seen/Unseen 混用。CW 配对分析亦已保存到附带 CSV，三数据集全体净变化同为4/$-7$/59张。
\section{本周产生的问题与数据分析结论}
\subsection{当前能够确认的结论}
\begin{enumerate}
\item 结果记录已经补齐。本轮完成 UT 的四变体和另外两个数据集的 A/C，主评价16组；后处理20次。MIT/C-GQA 的 D/E、全部新方法的 test、多 seed 和 EKR/USG 训练均未完成，也未当作结果填写。
\item 当前模板先验 $K$ 在 UT 的正强度使用没有主指标收益。重训 D/E 与固定权重后处理均呈现不利方向，应冻结负结果，而不反复更换选择标准找正值。
\item 文本版本 C 在 UT 获得明显单 seed validation增益；C-GQA幅度较小，MIT表现混合。后续优先复查可重复性和效果来源，不能直接宣称通用知识增强已成立。
\end{enumerate}
\subsection{尚不能确认的问题与原因}
\begin{enumerate}
\item \textbf{内容与归一化尚未分离。}A 使用冻结的基线名称原始均值向量，C 使用冻结的名称/描述单位向量聚合；两者同时改变语义内容与向量范数。虽最终 cosine 会归一化，前面的组合器和文本引导操作仍可能受输入尺度影响。现有 C-A 不能只归因于描述语义，尚需 name-only unit 对照及描述打乱对照。
\item \textbf{知识审核不足。}自动首义项、非视觉定义、通用模板和标签模板条目可能稀释有效描述。OOV约0.26\%--0.36\%只能说明词表覆盖较高，不保证多义词对应图像含义正确。
\item \textbf{单 seed 与验证集选择。}checkpoint、$\lambda$ 和路线选择均使用 validation，存在选择偏差；C-GQA的小差值尤其需要多个配对 seed。不能把单 run 的逐图数量当作训练稳定性的替代。
\item \textbf{基线有效配置归档不完整。}部分历史 YAML 没有命令行覆盖项，幸有日志可恢复实际设置。后续保存 merged training config，与评价 config分别冻结并写入哈希。
\item \textbf{环境故障曾中断诊断。}曾出现 torch/urllib3文件可读却无法通过目录枚举导入的异常，以及系统预加载库记录。保留事故日志和恢复后环境快照；本轮后续运行完成证明所需流程已执行，但不证明此前系统异常根因已由管理员确认。该问题与模型性能负结果分别记录。
\end{enumerate}
\section{下一步可行路线与实验安排}
\subsection{先确认现有文本收益来源}
冻结本周所有原始 run、缓存与评估协议。新增“仅名称单位向量”对照 $A_{unit}$：复用相同冻结名称、别名与 L2规则，保持其他设置与 C一致；对照现有 A、$A_{unit}$、C，分开观察向量尺度变化和描述内容的增量。描述打乱使用固定置换，保留条数、维度和 OOV处理，并独立命名；它是负对照，不是随机语义图方法。

先在 UT 用一个现有 seed检查接口，随后将 A/$A_{unit}$/C扩展为至少3个预先固定、配对的 seed（例如0、3351、2026）；已有 run只在 seed和所有协议一致时复用。报告各 seed结果、配对差值、均值和样本标准差，三个数据集分别统计。C-GQA需预留每个 run约半小时的最终 OW评价；MIT/C-GQA若延长预算，A/C同步延长并另立协议版本。

主要指标仍为 val CW AUC；同时报告 OW AUC、HM和未见指标，避免只挑涨点项。若 C持续优于 $A_{unit}$，且正确描述优于打乱描述，才进一步支持语义内容有效；若仅范数对照恢复收益，则调整研究定位，而不把归一化包装成知识推理创新。
\subsection{再验证 EKR 条件知识增量选择}
EKR 仍处于方案阶段，目标是利用被平均消去的描述差异。保留 C 的评分与候选协议，冻结 C第一阶段参数，对每个属性的名称和描述向量先中心化，选择由图像和候选对象条件决定的描述增量，并加入值为零的空选项：
\[
\mu_a=\frac1{R_a}\sum_r t_{ar},\qquad d_{ar}=t_{ar}-\mu_a,
\]
\[
s_{EKR}(x,a,o)=s_C(x,a,o)+\eta\sum_{r=0}^{R_a}
\alpha_r(x,o,a)\,\frac{u(x)^\top d_{ar}}{2},\quad d_{a0}=0.
\]
要求零增量恢复 C，均匀选择真实中心化描述时增量为零。第一版不加入固定 $K$ 偏置；图像与候选对象只是路由条件，不将其冒称为已经实现新的完整条件概率分解。

最小对照为 C、等参数普通适配器、普通描述注意力、完整 EKR，再做去对象条件/去中心化/去空选项和描述置换。所有结果都需与同一冻结 C、同训练预算比较；只有 EKR优于强对照，且收益依赖正确描述、在配对 seed间复现，才支持相应机制贡献。若没有超过 C或普通适配器，停止增加复杂度，先修订知识质量。
\subsection{与随机语义图路线严格分开}
USG学习 OI/OE/AE关系的边保留概率，涉及图采样、消息传播和跨图平均；EKR选择描述差异，不构建邻接矩阵，不采样边，也不做图分布 MC。当前 OI文件存在不等于语义图方法已经运行。若后续启动 USG，应独立目录、独立假设和确定图/EdgeDrop/学习均值图等对照，不能混入本周文本增量结果。

新方案没有事先可靠涨点保证。本周结果提供的可行方向是：保留 UT文本收益作为候选起点，先确认归一化与描述语义的各自作用，再以能退回 C的小分支检验额外证据。全部路线只用 train/val决定版本；协议冻结后再评价 test。
\section{数据归档与报告复核记录}
\subsection{归档来源与保存范围}
本周报告使用本地 \file{CPF_all_results} 中唯一完整导出，标识为：

\par\vspace{0.5\baselineskip}\noindent\begin{minipage}{\linewidth}

\begin{lstlisting}

czsl_all_completed_20261005_000428_2772161

\end{lstlisting}\end{minipage}\par\smallskip

导出完成于2026.10.5 00:20:52；本地下载和解压完成标记均为 true，下载清单已校验。归档含6类压缩包、2548个源文件，压缩总量约10.00 GB、展开文件约11.32 GB（十进制）。保存历史与本轮权重、日志、配置、指标、源码、知识源、缓存及数据划分元信息；不包含原始图像、原始 GloVe大文件和完整 Conda安装。它可追溯实验，不单独构成无需数据输入的完整重训环境。

实验身份使用 dataset、seed、variant、phase、setting、score mode、$\lambda$ 和 best epoch共同确定。\file{metrics.json} 保存原始精度、候选/图像数量及哈希。有效评价参数保存于 \file{config_eval_effective.yml}。源码快照、数据清单、知识 manifest 和 checkpoint均继续保留，新增实验不得覆盖本轮。
\subsection{主要权重与证据索引}

\begin{reporttable}

\begin{tabularx}{\linewidth}{|Y|r|c|}
\hline

\textbf{run} & \textbf{最佳轮} & \textbf{checkpoint SHA256前12位} \\ \hline

UT-Zappos/A & 6 & \file{9e8300061be1} \\ \hline

UT-Zappos/C & 6 & \file{7a87ed1e0bcd} \\ \hline

UT-Zappos/D & 6 & \file{4c7f251fe428} \\ \hline

UT-Zappos/E & 6 & \file{6322ca225e32} \\ \hline

MIT-States/A & 10 & \file{3c468979a12b} \\ \hline

MIT-States/C & 10 & \file{2ef05a49b292} \\ \hline

C-GQA/A & 9 & \file{0d1c5365a5f5} \\ \hline

C-GQA/C & 10 & \file{3231dcfa3998} \\ \hline

\end{tabularx}\end{reporttable}

\noindent{\small 表内哈希只用于阅读定位，不能将截断值当作完整校验值。完整64位值及实际重新哈希结果见 \file{actual_checkpoint_hashes.csv}。\par}

\begin{reporttable}

\begin{tabularx}{\linewidth}{|p{6.5cm}|Y|}
\hline

\textbf{分析材料} & \textbf{内容} \\ \hline

\file{all_val_metrics.csv} & 16行主评价，含最佳HM对应S/U、配置/来源/权重哈希 \\ \hline

\file{ac_comparison.csv} & 6组匹配A/C的完整差值 \\ \hline

\file{ut_acde_effects.csv} & UT各评价的文本/校准/交互差值 \\ \hline

\file{lambda_analysis.csv} & 10行候选汇总、入选条件、失败约束；对应20份原始JSON \\ \hline

\file{epoch_curves.csv} & 8个run×10轮日志曲线及最近打印loss \\ \hline

\file{knowledge_schema.json} & 来源字典、数量、字段示例及冻结哈希核验 \\ \hline

\file{prior_train_val_distribution.csv} & 仅train/val逐组合先验分布 \\ \hline

\file{bias0_cw/}、\file{bias0_ow/} & 逐图配对预测、按属性汇总和纠正/退化计数 \\ \hline

\end{tabularx}\end{reporttable}

\end{document}
